\section{Related Work}
\label{sec:related_work}

Our proposed framework lies at the intersection of weight-space model merging, adaptive ensembling, and statistical learning theory. In this section, we situate our work in the context of these three research lines.

\subsection{Weight-Space Model Merging}
Model merging has emerged as a compelling paradigm for multi-task ensembling without the inference cost of concurrent serving. Model Soups \cite{wortsman2022model} showed that averaging the parameters of multiple fine-tuned models from the same pre-trained initialization improves generalization. Task Arithmetic \cite{ilharco2022editing} extended this concept by adding or subtracting "task vectors" (differences between fine-tuned and pre-trained weights) to incorporate new capabilities into a base model. However, directly summing task vectors often results in severe parameter interference, especially in deep layers. 

To resolve coordinate-level conflicts, Yadav et al.\ proposed TIES-Merging \cite{yadav2023ties}, which prunes small-magnitude weight changes, resolves sign disagreements, and scales the remaining parameters. Similarly, Sparse Task Arithmetic \cite{drago2024sparse} analyzed these coordinate collisions and proved that sign consensus is often redundant due to the sparsity of high-magnitude updates, framing weight pruning as a magnitude-based noise filter. Despite their success, these weight-space heuristics rely on static, globally uniform merging coefficients. They fail to adapt to the highly varied functional requirements of different layers and tasks, motivating the shift toward adaptive model ensembling.

\subsection{Adaptive Test-Time Adaptation and Few-Shot Adaptation}
To address the limitations of static merging coefficients, several adaptive ensembling methods have been proposed. AdaMerging \cite{yang2023adamerging} introduced automated test-time adaptation (TTA) by optimizing layer-wise ensembling coefficients on unlabelled test streams via prediction entropy minimization. While conceptually elegant, unconstrained online TTA has been shown to suffer from significant instability. Vance et al.\ \cite{suitemerge2025} conducted a critical audit of test-time adaptation (SuiteMerge) and revealed that unsupervised entropy minimization is highly vulnerable to \emph{transductive overfitting} on local stream noise, frequently collapsing into degenerate states (where the model predicts a single class to trivially minimize entropy).

To regularize test-time coefficients, Croft \& Vance \cite{polymerge2024} proposed PolyMerge, which restricts online coefficients to low-degree polynomial trajectories. However, without labeled supervision, online methods remain highly susceptible to class collapse. Offline Few-Shot Validation Tuning (OFS-Tune) \cite{suitemerge2025} addresses this by optimizing coefficients on a small, labeled calibration dataset. Nevertheless, OFS-Tune optimizes a large number of parameters ($K \times L$) on extremely small datasets ($M \approx 10$ samples), creating a massive generalization gap due to overparameterization. Our work, Rademacher-Bounded Polynomial Merging, provides the first formal mathematical justification and generalization bounds for why polynomial trajectory constraints, combined with statistical regularization, successfully resolve this overparameterization.

\subsection{Generalization Bounds and Rademacher Complexity}
Evaluating the generalization capability of deep neural networks using learning-theoretic bounds is a classical problem in statistical learning theory \cite{vapnik1999overview, shalev2014understanding}. Traditional VC-dimension bounds are often vacuous for modern overparameterized models, leading researchers to focus on capacity-control measures such as Rademacher complexity \cite{bartlett2002rademacher, koltchinskii2002empirical}. Standard works have developed norm-based generalization bounds for deep networks \cite{neyshabur2015norm, bartlett2017spectrally} that scale with the product of the layers' spectral norms rather than the raw number of parameters, demonstrating that regularizing parameter trajectories and weight norms suppresses overparameterization.

Our work leverages these classical statistical learning principles to analyze the hypothesis space of merging coefficients. Rather than bounding the network's weights directly, we formulate the space of ensembling trajectories as a structured hypothesis class and bound its empirical Rademacher complexity. To our knowledge, this is the first work to apply Rademacher complexity bounds to weight-space model merging, establishing a rigorous connection between statistical learning theory and parameter-space ensembling.